% TOP_PROBLEM: {"id":30007643,"problem_number":"LOCAL-30007643","title":"The true global wealth distribution","category_id":21,"category":{"id":21,"name":"economics","display_name":"Economics","description":"Open economic theory statements, published research agendas, and proposed scoped specifications; question provenance and historical priority are qualified per record.","slug":"economics","order_index":21,"created_at":"2026-10-08T00:00:00Z"},"status":"research_question","external_url":null,"legacy_ids":[],"published":false,"statement":"In the explicitly specified setting: How concentrated is global wealth once offshore assets, private businesses, trusts, and underrepresented rich households are consistently measured? The mathematical statement above fixes the target, assumptions and scope of this catalogue formulation.","economics":{"original_id":132,"field":"Inequality and economic mobility","kind":"identification","formalization_basis":"adapted_research_question","source_status":"research_agenda","priority_status":"earliest_found","priority_note":"Earliest directly inspected qualifying passage in this audit; earlier origins were not established. A review or research-agenda statement does not imply original priority.","earliest_verified_reference":{"authors":["Gabriel Zucman"],"title":"Global Wealth Inequality","year":2019,"url":"https://gabriel-zucman.eu/files/Zucman2019.pdf","doi":null,"locator":"Section 6, Conclusion, printed p. 134","evidence_summary":"The conclusion calls for a global view including domestic and foreign wealth because conventional tax and survey sources are insufficient."},"current_reference":{"authors":["Gabriel Zucman"],"title":"Global Wealth Inequality","year":2019,"url":"https://gabriel-zucman.eu/files/Zucman2019.pdf","doi":null,"locator":"Section 6, Conclusion, printed p. 134","evidence_summary":"The conclusion calls for a global view including domestic and foreign wealth because conventional tax and survey sources are insufficient."},"formalization_note":"The cited passage establishes a research direction, not this exact mathematical statement. The notation, population, policy menu, assumptions, and estimands are an original adaptation for this catalogue; no universal unsolved theorem or source priority is claimed."},"statusQualification":"Published question or research agenda verified at the cited locator. The mathematical specification is adapted for this catalogue and includes additional assumptions; source publication does not establish first-ever priority or certify this exact model remains unsolved. The cited passage establishes a research direction, not this exact mathematical statement. The notation, population, policy menu, assumptions, and estimands are an original adaptation for this catalogue; no universal unsolved theorem or source priority is claimed.","statusReviewedAt":"2026-10-08"}
% FIRST_POSTED: 2026-10-08
% ENTRY_KIND: statement_only
% !TeX program = lualatex
\documentclass[11pt]{article}
\usepackage{fontspec}
\usepackage[a4paper,margin=25mm]{geometry}
\usepackage{amsmath,amssymb,mathtools}
\usepackage{xurl}
\usepackage[hidelinks]{hyperref}
\newcommand{\sref}[2]{\hyperlink{source:#1:#2}{[S#2]}}
\setlength{\emergencystretch}{3em}
\begin{document}
% problemId:30007643
\section*{The true global wealth distribution}\label{30007643}
\subsection{Definitions and mathematical statement}
Fix a date and define the global population as all resident adults, assigning each person to one country of residence. For person \(i\), let \(A_i\) be the market value of assets beneficially owned worldwide and \(D_i\) debts; net wealth is \(W_i=A_i-D_i\). Beneficial ownership rules must specify trusts, foundations, private businesses, and jointly owned property. Let \(F_W\) be the resulting world distribution after converting values with a declared exchange-rate or purchasing-power rule. Given household surveys, registries, estate records, international financial positions, and enforcement disclosures, recover the top-share function \(S_p=\sum_{i:W_i\ge q_{1-p}}W_i/\sum_iW_i\), where \(q_{1-p}\) is the appropriate weighted quantile and the denominator is positive. Specify rules for quantile ties. Construct identified intervals allowing offshore assets, rich-household nonresponse, uncertain private-business values, and missing ownership links. Reconcile domestic and foreign asset aggregates without double-counting custodial holdings or corporations. Determine which additional measurements narrow the intervals most. The target is a date-specific distribution with quantified uncertainty, not a proof that an administrative total uniquely identifies ownership. Earlier global estimates are substantive evidence but do not eliminate the remaining coverage and valuation problem.

\par\textbf{Formulation and scope.} The cited passage establishes a research direction, not this exact mathematical statement. The notation, population, policy menu, assumptions, and estimands are an original adaptation for this catalogue; no universal unsolved theorem or source priority is claimed.

\par\textbf{Open-question provenance.} An explicit question or research agenda was verified at the source locator. This does not by itself certify that every later version remains unresolved \sref{30007643}{1}.

\par\textbf{Historical priority.} Earliest directly inspected qualifying passage in this audit; earlier origins were not established. A review or research-agenda statement does not imply original priority.

\subsection{Short English statement}
In the explicitly specified setting: How concentrated is global wealth once offshore assets, private businesses, trusts, and underrepresented rich households are consistently measured? The mathematical statement above fixes the target, assumptions and scope of this catalogue formulation.

\subsection{Sources}
\begin{itemize}\raggedright
\item[\textnormal{[S1]}]\hypertarget{source:30007643:1}{} Gabriel Zucman. 2019. Global Wealth Inequality. Section 6, Conclusion, printed p. 134.
\url{https://gabriel-zucman.eu/files/Zucman2019.pdf}.
{\small\textit{Earliest verified question/agenda reference; historical priority qualified below. The conclusion calls for a global view including domestic and foreign wealth because conventional tax and survey sources are insufficient.}}
\end{itemize}
\end{document}
